\documentclass{beamer}

\usetheme{default}
\setbeamertemplate{footline}[frame number]
\usenavigationsymbolstemplate{}

\usepackage[utf8]{inputenc}
\usepackage[croatian]{babel}

%% drawing various dependency trees and graphs
\usepackage{tikz-dependency}
\usepackage{tkz-graph}

%% rollin-rollout graphs
\usepackage{tikz}
\usetikzlibrary{arrows.meta,decorations.pathreplacing,matrix,positioning}
\colorlet{lightblue}{blue!60!white}
\colorlet{darkgreen}{green!60!black}

%% subfigure and subcaptions
\usepackage{subcaption}

\usepackage{multirow} %% tablica
\usepackage{colortbl}
\usepackage{xcolor}
\usepackage{textcomp}
%% bolds the math consistently (italic, functions)
\newcommand{\mbold}[1]{{\boldmath #1}}

%% l2s
\def \lts {\textsc{l2s}}

%% for notes
\newcommand{\bq}[2]{{\fbox{#1}\colorbox{red}{~#2}}}

%% engl
\newcommand{\engla}[2]{(engl.~\emph{#1}, \textsc{#2})}

%% for exponents to differantiate outputs inputs and hypotheses
%% x\ssymbol{1} called like this
\def\@fnsymbol#1{\ensuremath{\ifcase#1\or *\or \dagger\or \ddagger\or
   \mathsection\or \mathparagraph\or \|\or **\or \dagger\dagger
   \or \ddagger\ddagger \else\@ctrerr\fi}}

\newcommand{\ssymbol}[1]{^{\@fnsymbol{#1}}}


%% argmin argmax operator
\DeclareMathOperator*{\argmax}{argmax}
\DeclareMathOperator*{\argmin}{argmin}

%% condition
\newtheorem{condition}{Uvjet}


%% fields
\newcommand*{\field}[1]{\mathbb{#1}}
%% max margin markov networks abbrev
\newcommand*{\mmmm}[0]{$\text{M}^4$}
%% text underscript
\newcommand{\textunderscript}[2]{$\text{#1}_{\text{#2}}$}

%% url font style
\urlstyle{sf}

%% code listings
\usepackage{listings}

%% algorithm

\usepackage{algorithm}
\usepackage{algpseudocode}

% algorithm counter resets every chapter
\makeatletter
\@addtoreset{algorithm}{chapter}
\makeatother
\renewcommand{\thealgorithm}{\arabic{algorithm}}
% Algorithm # is <chapter>.<algorithm>

\algrenewcommand{\algorithmicrequire}{\textbf{Potrebno:}}
\algrenewcommand{\algorithmicend}{\textbf{kraj}}
\algrenewcommand{\algorithmicif}{\textbf{ako}}
\algrenewcommand{\algorithmicthen}{\textbf{onda}}
\algrenewcommand{\algorithmicelse}{\textbf{inače}}
\algrenewcommand{\algorithmicfor}{\textbf{za}}
\algrenewcommand{\algorithmicforall}{\textbf{za svaki}}
\algrenewcommand{\algorithmicforall}{\textbf{za svaki}}
\algrenewcommand{\algorithmicfunction}{\textbf{funkcija}}
\algrenewcommand{\algorithmicdo}{}
\algrenewcommand{\algorithmicwhile}{\textbf{dok je}}
\algrenewtext{EndFor}{\textbf{kraj petlje}}
\algrenewtext{EndWhile}{\textbf{kraj petlje}}
\algrenewtext{EndIf}{\textbf{kraj}}
\algrenewtext{EndFunction}{\textbf{kraj funkcije}}
\algrenewcommand{\algorithmicreturn}{\textbf{vrati}}

\algdef{SE}[WHILE]{While}{EndWhile}[1]
  {\algorithmicwhile\ #1\ \textbf{ponavljaj}}
  {\textbf{kraj petlje}}

\makeatletter
\renewcommand{\ALG@name}{Algoritam}
\makeatother

\makeatletter
\renewcommand{\ALG@beginalgorithmic}{\scriptsize}
\algrenewcommand\alglinenumber[1]{\tiny #1:}
\makeatother

\title{Učenje pretraživanja za rješavanje zadataka obrade prirodnog jezika}

\author{Vjeran Crnjak\\Mentor: doc.~dr.~sc.~Jan Šnajder}

\institute{Fakultet elektrotehnike i računarstva, Sveučilište u Zagrebu\\
\includegraphics[scale=0.05]{images/fer-logo.pdf}
\qquad
\includegraphics[scale=0.1175]{images/unizg-logo.pdf}}
% - Use the \inst command only if there are several affiliations.
% - Keep it simple, no one is interested in your street address.

\date{\today}
% - Either use conference name or its abbreviation.
% - Not really informative to the audience, more for people (including
%   yourself) who are reading the slides online

\subject{Theoretical Computer Science}
% This is only inserted into the PDF information catalog. Can be left
% out.

% If you have a file called "university-logo-filename.xxx", where xxx
% is a graphic format that can be processed by latex or pdflatex,
% resp., then you can add a logo as follows:

% \pgfdeclareimage[height=0.5cm]{university-logo}{university-logo-filename}
% \logo{\pgfuseimage{university-logo}}

% Delete this, if you do not want the table of contents to pop up at
% the beginning of each subsection:
\AtBeginSubsection[]
{
  \begin{frame}<beamer>{Outline}
    \tableofcontents[currentsection,currentsubsection]
  \end{frame}
}

% Let's get started
\begin{document}

\begin{frame}
  \titlepage
\end{frame}

\begin{frame}{Outline}
  \tableofcontents
  % You might wish to add the option [pausesections]
\end{frame}

% Section and subsections will appear in the presentation overview
% and table of contents.
\section{First Main Section}

\subsection{First Subsection}

\begin{frame}{First Slide Title}{Optional Subtitle}
  \begin{itemize}
  \item {
    My first point.
  }
  \item {
    My second point.
  }
  \end{itemize}
\end{frame}

\subsection{Second Subsection}

% You can reveal the parts of a slide one at a time
% with the \pause command:
\begin{frame}{Second Slide Title}
  \begin{itemize}
  \item {
    First item.
    \pause % The slide will pause after showing the first item
  }
  \item {
    Second item.
  }
  % You can also specify when the content should appear
  % by using <n->:
  \item<3-> {
    Third item.
  }
  \item<4-> {
    Fourth item.
  }
  % or you can use the \uncover command to reveal general
  % content (not just \items):
  \item<5-> {
    Fifth item. \uncover<6->{Extra text in the fifth item.}
  }
  \end{itemize}
\end{frame}

\section{Second Main Section}

\subsection{Another Subsection}

\begin{frame}{Blocks}
\begin{block}{Block Title}
You can also highlight sections of your presentation in a block, with it's own title
\end{block}
\begin{theorem}
There are separate environments for theorems, examples, definitions and proofs.
\end{theorem}
\begin{example}
Here is an example of an example block.
\end{example}
\end{frame}

% Placing a * after \section means it will not show in the
% outline or table of contents.
\section*{Summary}

\begin{frame}{Summary}
  \begin{itemize}
  \item
    The \alert{first main message} of your talk in one or two lines.
  \item
    The \alert{second main message} of your talk in one or two lines.
  \item
    Perhaps a \alert{third message}, but not more than that.
  \end{itemize}

  \begin{itemize}
  \item
    Outlook
    \begin{itemize}
    \item
      Something you haven't solved.
    \item
      Something else you haven't solved.
    \end{itemize}
  \end{itemize}
\end{frame}



% All of the following is optional and typically not needed.
\appendix
\section<presentation>*{\appendixname}
\subsection<presentation>*{For Further Reading}

\begin{frame}[allowframebreaks]
  \frametitle<presentation>{For Further Reading}

  \begin{thebibliography}{10}

  \beamertemplatebookbibitems
  % Start with overview books.

  \bibitem{Author1990}
    A.~Author.
    \newblock {\em Handbook of Everything}.
    \newblock Some Press, 1990.


  \beamertemplatearticlebibitems
  % Followed by interesting articles. Keep the list short.

  \bibitem{Someone2000}
    S.~Someone.
    \newblock On this and that.
    \newblock {\em Journal of This and That}, 2(1):50--100,
    2000.
  \end{thebibliography}
\end{frame}

\begin{frame}{Lokalno optimalno učenje pretraživanja (\textsc{LOLS})}

\begin{center}
\begin{algorithm}[H]
\begin{algorithmic}[1]
\Require Skup podataka $\{x_i, y_i\}_{i=1}^N$ uzet iz distribucije $\mathcal{D}$
         i $\beta \geq 0$. %% -- parametar mješavine za rollout.
\State Inicijaliziraj policu $\hat{\pi}_1$.
\ForAll{$i \in \{1,2,\ldots,N\}$}
  \State Generiraj referentnu policu $\pi^{\text{ref}}$ temeljenu na $y_i$.
  \State Inicijaliziraj $\Gamma = \emptyset$. \Comment{skup primjera osjetljivih na trošak.}
  \ForAll{$t \in \{0,1,2,\ldots,T_i-1\}$}
    \State Primijeniti $t$ puta policu $\pi_{i}^{\text{in}} = \hat{\pi}_i$  i stići do $s_t$. \Comment{Rollin.} \label{alg:lols:learned}
    \ForAll{$a \in A(s_t)$}
      \State Neka je  $\pi_{i}^{\text{out}} = \pi^{\text{ref}}$ s vjerojatnošću $\beta$, inače $\hat{\pi}_i$.
      \State Procjeni trošak $c_{i,t}(a)$ koristeći $T-t-1$ puta policu $\pi_{i}^{\text{out}}$. \Comment{Rollout.} \label{alg:lols:mixture}
    \EndFor
    \State Generiraj vektor značajki $\mathbf{\Phi}(x_i, s_t)$.
    \State Postavi $\Gamma = \Gamma \cup \{\langle c_{i,t}, \mathbf{\Phi}(x_i, s_t) \rangle\}$.
  \EndFor
  \State $\hat{\pi}_{i+1} \gets \textsc{Nauči}(\hat{\pi}_i, \Gamma)$.
\EndFor
\State Vrati usrednjenu policu preko svih $\hat{\pi}_1, \hat{\pi}_2, \ldots, \hat{\pi}_N$.
\end{algorithmic}
\end{algorithm}
\end{center}

\end{frame}

\end{document}
